\documentclass[12pt]{article}
\usepackage[utf8]{inputenc}
\usepackage{times}
\usepackage{geometry}
\geometry{a4paper, margin=1in}
\usepackage{graphicx}
\usepackage{lipsum} 
\usepackage{cite} 
\usepackage{hyperref}
\usepackage{titling}
\usepackage{titlesec}
\usepackage{amsmath} 
\usepackage{array}

\title{\textbf{Your Thesis Title}}
\author{Your Name \\ Supervisor: Your supervisor's Name}
\date{\today}

\begin{document}

\maketitle

\begin{abstract}
This document is a template for writing a PhD proposal in any discipline. This abstract should summarize the motivation, main contributions, and results of your thesis in a concise format.
\end{abstract}

\tableofcontents
\newpage

\section{Introduction}

\subsection{Background}
Discuss the broader context of your research, its significance, and its contribution to your field.\cite{Turing1950}. 

\subsection{Problem Statement}
Define the specific problem or challenge your thesis addresses.

\subsection{Research Objectives}
List the objectives or aims of your research.
\begin{itemize}
    \item Objective 1
    \item Objective 2
    \begin{itemize}
        \item Sub-objective 1
        \item Sub-objective 2
        \item Sub-objective 3
    \end{itemize}
\end{itemize}

\section{Literature Overview}
Review the existing literature related to your research topic, highlighting how your work will contribute to the existing knowledge base.

\section{Method}
Explain the methods and approaches you are going to use in your research. A flowchart can help illustrate your method.

\begin{figure}[htbp]
    \centering
    \includegraphics[width=0.6\linewidth]{images/flowchart.png}
    \caption{A workflow for your research.}
    \label{fig:workflow}
\end{figure}

\section{Current Work and Initial Results}
Present the findings of your research. Use figures, tables, and charts as needed to illustrate your results.

\begin{table}[htbp]
    \centering
    \caption{Table Caption}
    \label{tab:table1}
    \begin{tabular}{lcc}
    \hline
    \textbf{Cell1} & \textbf{Cell2} & \textbf{Cell3} \\
    \hline
    Cell4 & Cell5 & Cell6 \\
    Cell7 & Cell8 & Cell9 \\
    \hline
    \end{tabular}
\end{table}

\section{Work Plan and Time Table}
The Work Plan and Time Table provide a structured outline and schedule for the project, detailing tasks, milestones, and their respective timelines.

\begin{table}[htbp]
    \centering
    \caption{Work plan}
    \label{tab:workplan}
    \begin{tabular}{|l|c|c|c|}
  \hline
  \textbf{Task} & \textbf{Start Date} & \textbf{End Date} & \textbf{Duration} \\
  \hline
  Task 1 & 01/01/2024 & 01/01/2025 & 1 year \\
  \hline
  Task 2 & 01/01/2025 & 01/01/2026 & 1 year \\
  \hline
  Task 3 & 01/01/2026 & 01/01/2027 & 1 year \\
  \hline
  \end{tabular}
\end{table}

\section{References}
\bibliographystyle{plain}
\bibliography{references.bib}

\end{document}
